% ============================================================
%  PID CONTROLLERS — COURSE SYLLABUS
% ============================================================
\documentclass[10pt]{article}
\usepackage[T1]{fontenc}
\usepackage[utf8]{inputenc}
\usepackage{lmodern}
\usepackage{microtype}
\usepackage{palatino}
\usepackage[scaled=0.82]{beramono}
\usepackage[top=0.55in, bottom=0.50in, left=0.65in, right=0.65in]{geometry}
\usepackage[dvipsnames,svgnames,x11names]{xcolor}
\usepackage{tikz}
\usepackage{enumitem}
\usepackage[most]{tcolorbox}
\usepackage[colorlinks=true, linkcolor=pblue, urlcolor=pcyan,
  pdftitle={PID Controllers --- Course Syllabus}]{hyperref}
\usepackage{parskip}

\definecolor{pblue}{HTML}{9333EA}
\definecolor{pcyan}{HTML}{A855F7}
\definecolor{pnavy}{HTML}{3B0764}
\definecolor{porange}{HTML}{D97706}
\definecolor{paccent}{HTML}{FAF5FF}
\definecolor{pgray}{HTML}{6B7280}

\setlength{\parskip}{2pt}
\setlength{\parindent}{0pt}
\pagestyle{empty}

\begin{document}

\begin{tikzpicture}[remember picture, overlay]
  \fill[pnavy] (current page.north west) rectangle ([yshift=-1.0in]current page.north east);
  \node[anchor=west, text=white, font=\fontsize{22}{26}\bfseries\sffamily]
    at ([xshift=0.65in, yshift=-0.35in]current page.north west) {PID Controllers};
  \node[anchor=west, text=pcyan!80, font=\fontsize{11}{13}\sffamily]
    at ([xshift=0.65in, yshift=-0.63in]current page.north west)
    {Study Guide · Industrial Automation Series · SCADA Hub Training Platform};
  \node[anchor=east, draw=pcyan, rounded corners=5pt, fill=pblue!70,
        inner sep=6pt, text=white, font=\small\bfseries\sffamily]
    at ([xshift=-0.65in, yshift=-0.35in]current page.north east)
    {10 Chapters · 350+ Questions · PDF Included};
  \draw[porange, line width=2pt]
    ([xshift=0.65in, yshift=-0.83in]current page.north west) --
    ([xshift=3.5in,  yshift=-0.83in]current page.north west);
\end{tikzpicture}

\vspace{0.95in}

{\small\sffamily\color{pgray}%
  10 Chapters\enspace·\enspace350+ Questions\enspace·\enspace3 Quiz Levels\enspace·\enspace Free Access\enspace·\enspace No Account Required}

\smallskip
{\color{pnavy}\rule{\linewidth}{1.2pt}}
\smallskip
{\small\sffamily\color{pgray}\MakeUppercase{Course Overview --- Chapter by Chapter}}
\smallskip

\setlist[enumerate]{noitemsep, topsep=2pt, leftmargin=*, label=\textcolor{pblue}{\small\bfseries\arabic*.}}
\begin{enumerate}\small
  \item \textbf{Introduction} --- Every process that requires maintaining a setpoint --- temperature, pressure, flow, level --- uses feedback control. This chapter establishes the closed-loop principle: measure, compare, correct. Engineers learn why PID is the default algorithm before understanding how it works.
  \item \textbf{Control Loop Fundamentals} --- Process variable, setpoint, error, and manipulated variable are the four elements every control loop shares regardless of the physical process. This chapter makes the vocabulary automatic before the math is introduced.
  \item \textbf{PID Algorithm} --- Proportional responds to current error. Integral eliminates steady-state offset. Derivative anticipates future error from its rate of change. Engineers learn what each term contributes to the output --- and what happens when one term dominates.
  \item \textbf{Tuning Methods} --- Ziegler-Nichols open- and closed-loop, lambda tuning, and IMC replace trial-and-error with a repeatable procedure. Engineers learn which method fits which process type and what to do when the textbook method produces an unstable loop.
  \item \textbf{Process Characterization} --- First-order plus dead time (FOPDT) model: gain, time constant, and dead time. These three parameters determine every tuning target. Engineers learn to identify them from a step test without simulation software.
  \item \textbf{Cascade Control} --- Inner and outer loop architecture for processes with slow primary dynamics and a fast measurable intermediate variable. Cascade is the single most effective structural improvement for difficult-to-tune loops.
  \item \textbf{Digital Implementation} --- Sampling rate, velocity form, bumpless transfer between manual and auto, and anti-windup. Real PLC PID blocks require these parameters; this chapter explains what they do and how to set them correctly.
  \item \textbf{PLC Implementation} --- Configuring PID function blocks in common PLCs: scaling inputs, setting output limits, switching modes, and the parameters vendors rename between products. What Rockwell calls Reset, Siemens calls Ti.
  \item \textbf{Troubleshooting} --- Oscillation, sluggish response, offset that integral should have removed, valve stiction, and the difference between a tuning problem and a mechanical problem. Field diagnosis without simulation.
  \item \textbf{Practice Lab} --- Four exercises: tuning a flow loop from a step test, designing a temperature cascade, configuring anti-windup on a PLC PID block, and commissioning a new loop from zero in manual before switching to auto.
\end{enumerate}

\smallskip
{\color{pblue}\rule{\linewidth}{0.6pt}}
\smallskip

\begin{minipage}[t]{0.48\textwidth}
\begin{tcolorbox}[
  enhanced, colback=paccent, colframe=pblue, arc=4pt,
  fonttitle=\small\bfseries\sffamily\color{white},
  title=Certification Relevance,
  boxed title style={colback=pblue},
  attach boxed title to top left={yshift=-2mm, xshift=5mm},
  top=6pt, bottom=4pt, left=5pt, right=5pt
]
\setlist[itemize]{noitemsep, topsep=0pt, leftmargin=*, label=\textcolor{pblue}{$\bullet$}}
\begin{itemize}\small
  \item \textbf{ISA CCST Level III} --- Advanced process control and loop tuning are tested at the senior technician level.
  \item \textbf{ISA CAP} --- Certified Automation Professional exam covers control system design and PID fundamentals explicitly.
  \item \textbf{PE Exam (Control Systems)} --- Professional Engineer licensing tests control theory, system response, and stability analysis at depth.
\end{itemize}
\end{tcolorbox}
\end{minipage}
\hfill
\begin{minipage}[t]{0.48\textwidth}
\begin{tcolorbox}[
  enhanced, colback=porange!8, colframe=porange, arc=4pt,
  fonttitle=\small\bfseries\sffamily\color{white},
  title=Next Steps,
  boxed title style={colback=porange},
  attach boxed title to top left={yshift=-2mm, xshift=5mm},
  top=6pt, bottom=4pt, left=5pt, right=5pt
]
\begin{enumerate}[noitemsep, topsep=0pt, leftmargin=*, label=\small\textcolor{porange}{\arabic*.}]
  \item \small Complete PID alongside the IEC 61131-3 guide --- PLC PID blocks make sense once you understand both the algorithm and the language.
  \item Download MATLAB (student license) or use Scilab (free) to simulate open-loop step tests before commissioning on live equipment.
  \item Review ISA-5.1 (Instrumentation Symbols) and ISA-18.2 (Alarm Management) as the surrounding standards.
  \item Pursue ISA CCST as the entry credential, then ISA CAP for the advanced process control track.
\end{enumerate}
\end{tcolorbox}
\end{minipage}

\begin{tikzpicture}[remember picture, overlay]
  \fill[pnavy!90] (current page.south west) rectangle ([yshift=0.30in]current page.south east);
  \node[anchor=west, text=white!70, font=\footnotesize\sffamily]
    at ([xshift=0.65in, yshift=0.15in]current page.south west)
    {SCADA Hub Training Platform · scada-hub.vercel.app · Free access · No account required for study material};
  \node[anchor=east, text=pcyan, font=\footnotesize\bfseries\sffamily]
    at ([xshift=-0.65in, yshift=0.15in]current page.south east)
    {Course 4 of 8};
\end{tikzpicture}

\end{document}
